\section{Disegno sperimentale}
\label{sec:design}

Gli esperimenti non sono una cronologia. Prima di addestrare qualunque cosa, le
osservazioni della Sezione~\ref{sec:data} sono state trasformate in una lista
ordinata di ipotesi, ognuna delle quali cambia esattamente una variabile ed è
falsificabile da un singolo numero:

\begin{enumerate}\itemsep0pt
    \item un encoder--decoder standard trova già i filamenti principali, e dà una
          baseline con cui vale la pena confrontarsi;
    \item con meno di 500 osservazioni il vincolo attivo è l'encoder, quindi le
          feature ImageNet devono battere l'inizializzazione casuale;
    \item il Dice è cieco alla frammentazione, quindi una loss che non lo è --
          Tversky più clDice -- deve alzare la PQ dove il Dice non può;
    \item parte di ciò che resta è il \emph{dominio} H-alpha, non la loss;
    \item ciò che limita il punteggio è la \emph{scala}: a $512 \times 512$ un
          filamento largo sei pixel nativi ne diventa uno e mezzo.
\end{enumerate}

Le ipotesi 3 e 5 fanno previsioni opposte su dove viva l'errore residuo, e solo una
sopravvive al contatto con i dati. Le ipotesi successive sono state derivate allo
stesso modo, ciascuna dall'analisi dell'errore del run precedente.

Due regole hanno tenuto onesto lo studio. Un candidato viene confrontato con il
migliore precedente solo \emph{al proprio operating point}, ricercato da capo,
perché una loss o una risoluzione diversa spostano la calibrazione e riutilizzare
una vecchia soglia falsificherebbe una buona idea per il motivo sbagliato. E un
guadagno si crede solo se sopravvive alla \textbf{cross-selection}: l'operating
point viene scelto su metà dei file di validazione e valutato sull'altra metà.
Misurato così, su 169 voci di annotazione, \textbf{ogni guadagno sotto
$\approx 0{,}005$ è indistinguibile dal rumore di selezione} -- una soglia che ha
poi ucciso tre meccanismi i cui numeri aggregati sembravano positivi.

Prima di iniziare abbiamo anche scritto cosa avremmo esplorato e quanto sarebbe
costato, in modo che lo studio entrasse nel tempo a disposizione
(Tabella~\ref{tab:plan}).

\begin{table}[ht]
    \centering
    \caption{Il piano: cosa abbiamo deciso di esplorare e il costo di un run.}
    \label{tab:plan}
    \small\begin{tabular}{lll}
        \toprule
        Cosa esploriamo & Valori considerati & Costo \\
        \midrule
        risoluzione d'ingresso & 512, 1024, 1536, 2048 nativo & 10--35 min / run \\
        encoder                & U-Net da zero, ResNet18, ResNet34 & 10--35 min / run \\
        loss                   & BCE+Dice, +Tversky, +clDice, Dice per istanza, \texttt{pos\_weight} & 15--30 min / run \\
        augmentation           & simmetrie D4, jitter fotometrico, scale jitter & incluso sopra \\
        paradigma              & U-Net per pixel vs.\ detector YOLOv8 & 30 min -- 2,5 h / run \\
        post-processing        & soglia, chiusura, dim.\ minima, gate & 10 min, senza training \\
        \bottomrule
    \end{tabular}
\end{table}

